% from ega4-2-fr.tex
\oldpage[IV-2]{211}

clos, les $A_{\mathfrak q_j}$ le sont et on sait (\emph{loc. cit.}) que ces
anneaux sont de dimension 1, donc ce sont des anneaux de valuation discrète,
d'où (iii). Enfin, (iv) provient de l'hypothèse de récurrence. La démonstration
sera achevée si nous prouvons le

\begin{lemma}[7.6.4.1]
\label{IV.2.s7.6.4.1-fr}
(Zariski). --- Soient $A$ un anneau semi-local noethérien, $t$ un élément de
son radical vérifiant les conditions (i) à (iv) précédentes; alors
$\widehat A$ est réduit.
\end{lemma}

\begin{proof}
Posons pour simplifier $A'=\widehat A$, et considérons $t$ comme un élément de
$A'$; on a $A'/tA'=(A/tA)\otimes_AA'$, et comme $A'$ est un $A$-module plat,
il résulte de (3.3.1) que les idéaux premiers de $A'$ associés au $A'$-module
$A'/tA'$ sont les idéaux premiers $\mathfrak q'_{jh}$ tels que
$\mathfrak q'_{jh}$ soit au-dessus de $\mathfrak q_j$ et associé au
$k_j$-module $(A'/tA')\otimes_Ak_j$, où $k_j$ est le corps des fractions de
$A/\mathfrak q_j$. Or,
$\operatorname{Spec}((A'/tA')\otimes_Ak_j)$ est la fibre formelle de $A$ au
point $\mathfrak q_j$, ou aussi celle de $A/\mathfrak q_j$ au point
$\mathfrak q_j$ (point générique de
$\operatorname{Spec}(A/\mathfrak q_j)$). Comme en vertu de (iv) le complété
$A'/\mathfrak q_jA'$ de $A/\mathfrak q_j$ est réduit, il en est de même de la
fibre formelle de $A/\mathfrak q_j$ au point générique de
$\operatorname{Spec}(A/\mathfrak q_j)$; cette fibre formelle n'a donc pas de
cycle premier associé immergé et ses anneaux locaux aux points génériques
$\mathfrak q'_{jh}$ de ses composantes irréductibles sont des corps (3.2.1).
On voit donc par (3.3.3) et l'hypothèse (ii) que $A'/tA'$ n'a pas d'idéaux
premiers associés immergés, et d'autre part que
$\mathfrak q_jA'_{\mathfrak q'_{jh}}$ est l'idéal maximal de
$A'_{\mathfrak q'_{jh}}$ pour tout $h$; comme $A_{\mathfrak q_j}$ est un
anneau de valuation discrète, son idéal maximal
$\mathfrak q_jA_{\mathfrak q_j}$ est principal, donc l'idéal maximal de
l'anneau local noethérien $A'_{\mathfrak q'_{jh}}$ est principal, ce qui
entraîne que cet anneau est un anneau de valuation discrète (Bourbaki,
\emph{Alg. comm.}, chap. VI, § 3, n\textsuperscript{o} 6, prop. 9). Nous venons
ainsi de vérifier les hypothèses (i) à (iv) pour l'anneau local complet $A'$.
Il suffit donc de montrer que si $A$ est complet et vérifie les hypothèses (i)
à (iii) ((iv) l'étant automatiquement dans ce cas), alors $A$ est réduit. Or,
les hypothèses (i) et (ii) impliquent que, si
\[
\varphi:A\longrightarrow\prod_{j=1}^nA_{\mathfrak q_j}
\]
est l'homomorphisme canonique, on a
$t^mA=\varphi^{-1}(\varphi(t^mA))$ (3.4.9); comme l'image canonique de $tA$
dans l'anneau de valuation discrète $A_{\mathfrak q_j}$ est une puissance
$\mathfrak q_j^hA_{\mathfrak q_j}$ de l'idéal maximal, on voit que sur $A$ la
topologie $(tA)$-préadique est l'image réciproque par $\varphi$ de la topologie
produit sur $\prod_{j=1}^nA_{\mathfrak q_j}$. Or, comme $t$ appartient au
radical de $A$, la topologie $(tA)$-préadique est séparée, donc $\varphi$ est
injectif. Mais en vertu de l'hypothèse (iii), les $A_{\mathfrak q_j}$ sont
intègres, donc $\prod_{j=1}^nA_{\mathfrak q_j}$ est réduit, et a fortiori il en
est de même de $A$. C.Q.F.D.
\end{proof}
\end{proof}

\begin{corollary}[7.6.5]
\label{IV.2.s7.6.5-fr}
Soit $A$ un anneau semi-local noethérien vérifiant les conditions équivalentes
de (7.6.4); toute $A$-algèbre semi-locale essentiellement de type fini vérifie
aussi les conditions de (7.6.4).
\end{corollary}

\begin{proof}
La condition $a'')$ de (7.6.4) signifie en effet que $A$ est un
$\boldsymbol P$-anneau, où $\boldsymbol P(Z,k)$ est la propriété « $Z$ est
géométriquement réduit sur $k$ »; le corollaire est donc conséquence du
théorème général (7.4.4).
\end{proof}

\begin{corollary}[7.6.6]
\label{IV.2.s7.6.6-fr}
Soient $A$ un anneau semi-local noethérien intègre de dimension 1, $K$ son
corps des fractions. Pour que $A$ soit un anneau japonais, il faut et il suffit
que le complété $\widehat A$

